\section{Dominios y condiciones de las realizaciones}
\label{v:integracion}
La exclusión algebraica, su transporte por refinamiento y sus aplicaciones
estadísticas son resultados relacionados, con dominios distintos.
La identidad general de aniquilación
\eqref{v:fock:eq:transporte-aniquilacion} explicita el adjunto del transporte;
su especialización isométrica conserva la condición necesaria sobre la
aplicación entre los espacios de una partícula. Las ocupaciones de
equilibrio utilizan, además, una graduación y un Hamiltoniano compatibles
con la representación de la trayectoria.

La realización dinámica finita se desarrolla en la
sección~\ref{sec:hmtf2-cristal-temporal}. El operador de área, el
Hamiltoniano modular y la purificación sectorial se desarrollan en la
sección~\ref{v:ant:sec:area} y sus continuaciones. Las secciones
\ref{v:ant:sec:action} y \ref{v:ant:sec:ii-kb-aut-desarrollo} contienen
las construcciones de acción y de la pareja Boltzmann--temperatura que
utiliza la realización radiativa. La caracterización variacional de la
medida de historias y la convergencia de las sumas espectrales enlazan,
respectivamente, incidencia con entropía y modos discretos con radiación.
La sección radiativa
presenta sus hipótesis de cavidad,
dispersión, polarización y equilibrio, así como las pruebas espectrales
que utiliza. La sección~\ref{v:sec:enlace-velocidad-constitutiva}
reúne además la identidad de las realizaciones cinemática y constitutiva
de la velocidad, su normalización dimensional y su conservación bajo
amplificación del espacio de dos hojas.

Cada composición conserva las hipótesis matemáticas y la realización
física que declara. En particular, los modos armónicos del complejo
celular, las conexiones módulo calibre y los modos radiativos mantienen
los dominios distinguidos en la
sección~\ref{v:rec:completaciones:section:08}; la realización radiativa
utiliza explícitamente su dispersión y sus polarizaciones.

La subsección~\ref{v:subsec:limite-termodinamico} demuestra, para los tres
observables \(F_N,F_E,F_Z\), el paso de las sumas periódicas tridimensionales
a las integrales radiales cuando \(L\to\infty\). Su dominio exige \(a,b>0\),
mantiene fijado el modo constante de fondo y controla uniformemente el origen
y la cola espectral. El resultado proporciona las densidades de número,
energía de excitación y logaritmo de la función de partición de los modos
transversales; es distinto del bloque finito de capacidad espectral.

La sección~\ref{v:sec:capacidad-espectral-termica} tiene por dominio el
portador graduado finito, los árboles microscópico y agregado, el
levantamiento reversible de memoria, el reloj de capacidad y las coordenatizaciones
termométricas allí declaradas. Sus identidades espectrales, de transporte y
variacionales se prueban dentro de esos dominios; el reconocimiento físico
interviene después. La coordenatización dimensional tipada conserva la procedencia de
la intensidad marcada, la coordenatización energética y la temperatura física.

La sección~\ref{v:n19:sec:incertidumbre-fourier} trabaja sobre la
realización finita \(\ell^2((\mathbb Z/27\mathbb Z)^2)\) y su
Fourier unitaria. Sus entropías son las de las distribuciones en dos bases
conjugadas; los lectores dodecafásicos poseen el dominio diferente
\(\{-1,0,+1\}^{12}\) y aportan un resultado complementario de memoria
ordenada. La sección~\ref{sec:f051-prime-vacancy-observer} establece la
identidad modal exacta para cada \(X\) finito; su testigo histórico alcanza
\(10^8\) y la reproducción independiente documentada alcanza \(10^6\).
La promoción a la cota crítica requiere estimaciones uniformes en los modos.
